% =======================================================
% 5. DESIGN PATTERN UTILIZZATI
% =======================================================
\section{Design Pattern Utilizzati}
Di seguito sono illustrati i design pattern adottati per risolvere problematiche ricorrenti e garantire la scalabilità del codice di \textit{Second Brain}.

\subsection{Pattern Creazionali}

\subsubsection{[Inserire Pattern Creazionale, es. Factory / Dependency Injection]}
[Spiegare come il pattern viene usato per istanziare dinamicamente le strategie operative o i client LLM in base alla configurazione].

\subsection{Pattern Strutturali}

\subsubsection{[Inserire Pattern Strutturale, es. Adapter]}
[Spiegare come il pattern viene utilizzato per mascherare le differenze tra le varie API esterne dei provider LLM, offrendo un'interfaccia unica al sistema].

\subsection{Pattern Comportamentali}

\subsubsection{[Inserire Pattern Comportamentale, es. Strategy]}
[Spiegare come il pattern isola la logica di costruzione dei prompt per le diverse operazioni (Traduzione, Analisi Sei Cappelli, Riassunto), permettendo di aggiungere nuove funzionalità senza modificare il flusso principale].

\subsubsection{[Inserire Pattern Comportamentale, es. State / Observer]}
[Spiegare come viene gestita la reattività della UI nel Frontend durante l'attesa di una risposta dall'LLM (es. stati di Idle, Loading, Success, Error) o la gestione degli eventi dell'editor Markdown].